%% kapa: exam=1 round=24 year=2013 period=1 form=A subject=civil_law subject_ko=민법 seq=22 number=22
%% source: Q-Net 감정평가사 기출문제 (cst00309) artlSeq=5073802 file=extracted/1008/5073802/2013년 제24회 감정평가사 1차시험 기출문제/1교시(A형).hwp
%% conversion: kordoc 4.24.51 -> tools/gen1.py (auto); review=proofread-sonnet; flags=none
\begin{question}{22}
등기의 추정력에 관한 설명으로 옳지 않은 것은? (다툼이 있으면 판례에 의함)
\begin{choices}
  \item 소유권에 관한 등기에는 권리의 추정력이 있으므로 이를 다투는 측에서 그 무효사유를 주장ㆍ증명하여야 한다.
  \item 등기절차가 적법하게 진행되지 아니한 것으로 볼만한 의심스러운 사정이 증명된 경우에는 등기의 추정력이 깨어진다.
  \item 매매를 원인으로 하는 소유권이전등기가 경료된 경우 등기 기재와 같이 매매에 의하여 등기명의인이 소유권을 취득한 것으로 추정된다.
  \item 권리취득의 원인을 등기부에 기재된 취득원인과 달리 주장하는 경우 그 주장만으로는 등기의 추정력이 깨어지지 않는다.
  \item 부동산의 등기명의자는 자신의 소유권이전등기가 유래한 전 소유자에 대하여 등기의 추정력을 주장할 수 없다.
\end{choices}
\end{question}
\answer{5}
